\documentclass[12pt]{article}
\usepackage[utf8]{inputenc}
\usepackage[T1]{fontenc}
\usepackage[margin=1in]{geometry}
\usepackage{setspace}
\usepackage{xcolor}
\usepackage[colorlinks=true, urlcolor=blue!60!black, linkcolor=blue!60!black]{hyperref}
\setlength{\parindent}{2em}
\setlength{\parskip}{1.25em}
\renewcommand{\baselinestretch}{1.0}
\begin{document}
\begin{center}
{\LARGE \textit{Ex parte Milligan}, 71 U.S. (4 Wall.) 2 (1866)\par}
\vspace{0.8em}
{\large Case Brief\par}
\end{center}
\vspace{1em}

\section{Memory Jogger}
An Indiana civilian sentenced to hang by a Civil War military commission won his release: the Court held unanimously that the commission lacked jurisdiction, and 5--4 that neither the President nor Congress may subject civilians to military trial where the civil courts are open, even while habeas corpus is suspended (\href{https://supreme.justia.com/cases/federal/us/71/2/}{Justia}).

\section{Detailed Case Facts}
In the Habeas Corpus Act of March 3, 1863, Congress authorized President Lincoln to suspend the privilege of the writ of habeas corpus during the Civil War, while requiring in \S\S~2--3 that federal officials furnish lists of non-prisoner-of-war citizens detained in loyal states where federal courts remained unimpaired, and directing that any such prisoner be discharged if a federal grand jury convened and adjourned without indicting him. Pursuant to the Act, President Lincoln issued a proclamation on September 15, 1863, suspending the privilege of the writ for spies, aiders of the enemy, and draft resisters (\href{https://supreme.justia.com/cases/federal/us/71/2/}{71 U.S. at 4--6}).

Lambdin P. Milligan was a U.S. citizen who had resided in Indiana for twenty years, had never served in the U.S. military or militia, and had never resided in a rebellious state. On October 5, 1864, military officers arrested Milligan at his Indiana home on the order of Brevet Major-General Alvin P. Hovey, commandant of the Military District of Indiana. On October 21, 1864, Milligan was tried before a military commission at Indianapolis on charges of conspiracy against the United States, affording aid and comfort to rebels, inciting insurrection, disloyal practices, and violating the laws of war---specifically, aiding the secret ``Sons of Liberty'' to seize arsenals and free Confederate prisoners. Over Milligan's jurisdictional objection, the commission convicted him and sentenced him to hang on May 19, 1865, with presidential approval (\href{https://supreme.justia.com/cases/federal/us/71/2/}{71 U.S. at 6--7, 107}).

Meanwhile, in January 1865, the U.S. Circuit Court for the District of Indiana convened peacefully in Indianapolis with an empaneled federal grand jury, which adjourned on January 27, 1865, without indicting Milligan.

\section{Procedural History}
On May 10, 1865, Milligan petitioned the Circuit Court for a writ of habeas corpus and discharge under the 1863 Act. Because the two Circuit Court judges divided on whether the writ should issue, whether Milligan should be discharged, and whether the military commission had jurisdiction, they certified those three questions to the Supreme Court under the Judiciary Act of 1802 (\href{https://supreme.justia.com/cases/federal/us/71/2/}{71 U.S. at 7--9, 108--09}).

\section{Judicial Votes}
The Supreme Court answered the certified questions in Milligan's favor, \textbf{9--0} on the judgment that the military commission lacked jurisdiction and Milligan was entitled to discharge, and \textbf{5--4} on whether Congress could have authorized such a commission. The case was argued March 5--9 and 12--13, 1866, and decided April 3, 1866, with full opinions delivered in December 1866 (\href{https://supreme.justia.com/cases/federal/us/71/2/}{Justia}; \href{https://www.law.cornell.edu/supremecourt/text/71/2}{Cornell LII}):

\begin{itemize}
    \item \textbf{Majority Opinion:} Justice Davis delivered the opinion of the Court, joined by Justices Nelson, Grier, Clifford, and Field.
    \item \textbf{Concurring Opinion:} Chief Justice Chase filed an opinion concurring in the judgment, joined by Justices Wayne, Swayne, and Miller.
    \item \textbf{Dissenting Opinions:} None.
\end{itemize}

\section{Holding}
\textbf{Question Presented:}

During wartime and even while the privilege of the writ of habeas corpus is lawfully suspended, may the President or Congress subject a civilian citizen residing in a loyal state---where regular federal civil courts are open and functioning unobstructed---to trial, conviction, and sentence by a military commission?

\textbf{Holding:} No. Unanimously as to judgment, the Supreme Court held that the Indiana military commission lacked legal jurisdiction to try and sentence Milligan and that he was entitled to habeas relief and discharge under the Habeas Corpus Act of 1863. By a 5--4 majority, the Court further held that \textit{neither the President nor Congress} may authorize military commissions to try civilians who are unconnected to the armed forces in states where the regular civil courts are open and operating unobstructed, even during wartime suspension of habeas corpus (\href{https://supreme.justia.com/cases/federal/us/71/2/}{71 U.S. at 121--27, 130--31}).

\section{Analysis of Opinions}
\subsection{Majority Opinion (Justice Davis)}

After confirming Supreme Court jurisdiction over the certificate of division and Circuit Court authority under the 1863 Act (\href{https://supreme.justia.com/cases/federal/us/71/2/}{71 U.S. at 109--18}), Justice Davis anchored the majority's constitutional holding in four interlocking principles:

\begin{enumerate}
  \item \textbf{The Constitution Governs in War as in Peace:} Rejecting the doctrine of wartime necessity, Davis declared that the Constitution binds rulers and citizens equally in war and peace. Neither the President, Congress, nor the Judiciary may suspend any constitutional safeguard of civil liberty except the single right expressly subject to suspension---the privilege of the writ of habeas corpus under Article I, \S~9, cl. 2 (\href{https://supreme.justia.com/cases/federal/us/71/2/}{71 U.S. at 120--21, 125}).

  \item \textbf{Article III and Jury-Trial Guarantees:} Every criminal trial exercises the judicial power of the United States, which Article III vests exclusively in the Supreme Court and congressionally ordained inferior courts staffed by life-tenured judges. Likewise, Article III, \S~2 and the Fourth, Fifth, and Sixth Amendments guarantee grand jury indictment and trial by an impartial jury to all civilians accused of federal crimes. The Fifth Amendment's narrow textual exception for ``cases arising in the land or naval forces, or in the militia, when in actual service'' confirms that civilians outside the military cannot be stripped of jury trials (\href{https://supreme.justia.com/cases/federal/us/71/2/}{71 U.S. at 119--23}).

  \item \textbf{The ``Open Courts'' Limit on Martial Law:} The government invoked the ``laws and usages of war'' and martial law because Indiana lay within a military district and faced threatened invasion. Davis drew a bright line: martial rule cannot rest on threatened invasion or executive discretion. It is permissible only on the theater of active military operations where actual invasion or civil war has physically closed the courts and deposed civil authority, and it terminates the moment the courts reopen. In Indiana, federal authority was undisputed and the Circuit Court was open and actively trying wartime offenses (\href{https://supreme.justia.com/cases/federal/us/71/2/}{71 U.S. at 121--22, 124--27}).

  \item \textbf{Limits of Habeas Suspension and Prisoner-of-War Status:} Suspending the privilege of habeas corpus permits temporary detention without judicial release, but ``[t]he Constitution goes no further''---it does not authorize military trial or dispense with the common-law criminal process. Nor did suspension abolish the writ itself; the writ issues as of course so a court may determine whether the petitioner falls within the suspension. Finally, a lifelong civilian resident of a loyal state who never bore arms could not be classified as a prisoner of war (\href{https://supreme.justia.com/cases/federal/us/71/2/}{71 U.S. at 125--26, 130--31}).
\end{enumerate}

\subsection{Concurring in the Judgment (Chief Justice Chase)}

No Justice dissented from the judgment. Chief Justice Chase, joined by Justices Wayne, Swayne, and Miller, concurred in the judgment but parted sharply with the majority's constitutional limitation on Congress (\href{https://supreme.justia.com/cases/federal/us/71/2/}{71 U.S. at 132--42}).

Chase agreed that Milligan had to be discharged because the Habeas Corpus Act of 1863 affirmatively required civilian grand-jury proceedings or release in loyal states, thereby prohibiting the President's unilateral military commission by clear statutory implication. However, Chase insisted the majority went unnecessarily far in holding that \textit{Congress} lacked constitutional power to authorize military commissions in states with open courts. Grounding legislative authority in Congress's Article I powers to declare war, raise and support armies, and govern the armed forces, Chase argued that during actual war or imminent public danger, civil courts might remain technically open yet be ``wholly incompetent'' or sympathetic to rebels. Distinguishing \textit{military law} (governing armed forces), \textit{military government} (in occupied territory), and \textit{martial law proper}, Chase contended that Congress---though not the President acting alone---possesses constitutional authority to establish military tribunals for wartime offenses threatening national security even where civilian courts sit (\href{https://supreme.justia.com/cases/federal/us/71/2/}{71 U.S. at 136--42}).

\subsection{Subsequent Developments}

\begin{itemize}

\item \textbf{Unlawful Belligerents (\textit{Ex parte Quirin}):} During World War II, in \href{https://supreme.justia.com/cases/federal/us/317/1/}{\textit{Ex parte Quirin}}, 317 U.S. 1 (1942), the Court unanimously upheld military-commission trials of eight German saboteurs (including one U.S. citizen) who covertly landed on U.S. shores out of uniform. Distinguishing \textit{Milligan}, \textit{Quirin} held that unlawful enemy belligerents charged with congressionally sanctioned law-of-war offenses are subject to military trial even when civilian courts are open, characterizing Milligan as a non-belligerent civilian never associated with enemy armed forces.

\item \textbf{Military Detention of Citizens (\textit{Hamdi v. Rumsfeld}):} In \href{https://supreme.justia.com/cases/federal/us/542/507/}{\textit{Hamdi v. Rumsfeld}}, 542 U.S. 507 (2004), a plurality held that Congress's 2001 Authorization for Use of Military Force permitted military \textit{detention} (with due process review) of a U.S. citizen captured on a foreign battlefield, invoking \textit{Quirin} to distinguish \textit{Milligan}; Justice Scalia's celebrated dissent (joined by Justice Stevens) relied squarely on \textit{Milligan} to argue that absent a Suspension Clause suspension, a citizen must be prosecuted for treason in Article III courts or released.

\end{itemize}

\section{Examples: Future Applications}
\begin{enumerate}

\item \textbf{Same-Side (Military Trial Invalid): Congress Authorizes Commissions for Domestic Saboteurs.} During a declared international conflict, Congress suspends the privilege of habeas corpus in a border state and enacts a statute authorizing military commissions to try U.S. citizens arrested domestically for ``aiding enemy propaganda and disrupting defense supply chains.'' A U.S. civilian living in that state---who has never joined enemy armed forces---is arrested at home and tried by a military commission while federal district courts in the state remain open and hearing criminal dockets. Under Justice Davis's 5--4 majority in \textit{Milligan}, the military trial is unconstitutional even with congressional authorization and habeas suspension. Suspension permits temporary detention, not military criminal adjudication; because the defendant is a non-belligerent civilian and Article III courts are open, the Fifth and Sixth Amendments mandate indictment and jury trial in civilian court. (Under Chief Justice Chase's concurrence, by contrast, Congress's express wartime authorization might save the commission.)

\item \textbf{Same-Side (Military Trial Invalid): Presidential Martial Law After Riots.} After several nights of rioting in a major city, the President, acting without any statute, declares martial law in the city and orders that persons arrested for arson be tried by military tribunals. A college student who has no connection to any armed force is arrested for setting a police car on fire and convicted by a tribunal, although the federal and state courts in the city never closed. The conviction cannot stand. Under Justice Davis's majority, martial rule ``can never exist where the courts are open,'' and the student is a civilian entitled to indictment and jury trial. Even Chief Justice Chase's concurrence would not save the tribunal, because Chase located the power to authorize military commissions in Congress, not the President acting alone (\href{https://supreme.justia.com/cases/federal/us/71/2/}{Justia}).

\item \textbf{Opposite-Side (Military Jurisdiction Permitted): A Citizen Saboteur Enlisted With the Enemy.} First, if a U.S. citizen secretly enlists in a hostile foreign military, slips across U.S. coastal defenses out of uniform during wartime to blow up a naval shipyard, and is charged with a recognized law-of-war violation under a statute authorizing military commissions, \href{https://supreme.justia.com/cases/federal/us/317/1/}{\textit{Ex parte Quirin}} permits trial by military commission despite open civil courts.

\item \textbf{Opposite-Side (Military Detention Permitted): A Citizen Captured on a Foreign Battlefield.} Second, if a U.S. citizen is captured bearing arms for enemy forces on an active foreign battlefield and Congress has authorized military force, \href{https://supreme.justia.com/cases/federal/us/542/507/}{\textit{Hamdi v. Rumsfeld}} permits preventive military \textit{detention} for the duration of hostilities (subject to due process review) rather than immediate civilian prosecution.

\item \textbf{Fence-Sitter (Open but Compromised Courts): A Militia Supporter in an Insurrection.} During an armed insurrection in one State, a militia seizes several county seats, and its members threaten jurors and witnesses in the federal courthouse, which stays open but cannot seat a jury in militia cases. Congress authorizes military commissions in that State for persons who knowingly supply the militia with weapons. A citizen who sold rifles to the militia, but never joined it, is tried by a commission. Under Davis's bright line, the courts are open, so the trial is invalid; but Chase warned that courts may be technically open yet ``wholly incompetent,'' and here Congress has acted. \textit{Quirin} also complicates matters, since the defendant's knowing support for an armed force may or may not make him a belligerent rather than a \textit{Milligan}-type civilian. The outcome is uncertain.

\end{enumerate}

\section{Critique}
\textbf{Post-War Courage Versus Wartime Realism:} Critics note that the Court issued \textit{Milligan} only after Appomattox, when ``the public safety [was] assured'' (\href{https://supreme.justia.com/cases/federal/us/71/2/}{71 U.S. at 109}), having declined to intervene in military commissions during active hostilities (\textit{Ex parte Vallandigham} (1864)). Chief Justice Chase's concurrence also highlights the practical flaw in a purely binary ``open courts'' test: courts may remain physically open while compromised by intimidation or disloyalty.

\textbf{The \textit{Quirin}--Davis Tension:} By carving out ``unlawful enemy belligerents'' even for U.S. citizens on domestic soil, \href{https://supreme.justia.com/cases/federal/us/317/1/}{\textit{Quirin}} created a persistent boundary problem between treasonous civilians protected by \textit{Milligan} and belligerent saboteurs triable by military commission.

\section{Key Quotations}
\begin{quote}
``The Constitution of the United States is a law for rulers and people, equally in war and in peace, and covers with the shield of its protection all classes of men, at all times and under all circumstances. No doctrine involving more pernicious consequences was ever invented by the wit of man than that any of its provisions can be suspended during any of the great exigencies of government.''
\end{quote}
\noindent --- Justice Davis, \href{https://supreme.justia.com/cases/federal/us/71/2/}{71 U.S. at 120--21}.

\begin{quote}
``Not one of these safeguards can the President or Congress or the Judiciary disturb, except the one concerning the writ of habeas corpus. ... The Constitution goes no further. It does not say, after a writ of habeas corpus is denied a citizen, that he shall be tried otherwise than by the course of the common law.''
\end{quote}
\noindent --- Justice Davis, \href{https://supreme.justia.com/cases/federal/us/71/2/}{71 U.S. at 125--26}.

\begin{quote}
``Martial rule can never exist where the courts are open and in the proper and unobstructed exercise of their jurisdiction. It is also confined to the locality of actual war.''
\end{quote}
\noindent --- Justice Davis, \href{https://supreme.justia.com/cases/federal/us/71/2/}{71 U.S. at 127}.

\begin{quote}
``We think that Congress had power, though not exercised, to authorize the military commission which was held in Indiana. ... The fact that the Federal courts were open was regarded by Congress as a sufficient reason for not exercising the power, but that fact could not deprive Congress of the right to exercise it.''
\end{quote}
\noindent --- Chief Justice Chase, concurring in the judgment, \href{https://supreme.justia.com/cases/federal/us/71/2/}{71 U.S. at 137, 140--41}.


\end{document}
